\documentclass[10pt,a4paper]{amsart}
\usepackage[margin=19mm]{geometry}
\usepackage{amsmath,amssymb,mathtools}
\usepackage[scr=rsfs,cal=euler]{mathalfa}
\usepackage{xfrac}
\usepackage[unicode,hidelinks,bookmarks=false]{hyperref}
\newcommand{\ie}{i.e.\ }
\newcommand{\cf}{cf.\ }
\newcommand{\resp}{respectively\ }
\newcommand{\Subsection}{\subsection}
\providecommand{\nobreakdash}{\nobreak\hbox{-}}
\setlength{\emergencystretch}{2em}
\multlinegap0pt
\usepackage{amssymb}
\usepackage{tikz}
\newtheorem{lem}{}
\title{The Hessian equation on nonsmooth $k$-convex domains or in the presence of subsolutions}
\author{J. Lukas Gehring}
\date{Source: arXiv:2607.21755v1 (CC BY 4.0)}
\begin{document}
\maketitle

\noindent\emph{Source and licence.} The Hessian equation on nonsmooth $k$-convex domains or in the presence of subsolutions. Version arXiv:2607.21755v1 (CC BY 4.0), distributed by arXiv under the Creative Commons Attribution 4.0 International licence, http://creativecommons.org/licenses/by/4.0/, as recorded in the arXiv licence metadata for this exact version and shown on the abstract page https://arxiv.org/abs/2607.21755v1. Full licence notice: \texttt{licenses/arxiv-2607.21755-CC-BY-4.0.txt}.\par\smallskip

\noindent\textit{Editorial scope.} Counted unit: the article's lem lem:Stone=000020Dini (source0.tex lines 807--819). The article's complete original proof of it is delivered with it (source0.tex lines 821--836, one \texttt{proof} environment of 16 lines). No mathematics is written, repaired or re-derived: the statement and the proof are the article's own lines, in the article's own notation, and no retained line is substituted or re-flowed.  No further source-local context is needed: every cross-reference of every retained line resolves inside the two delivered spans.  Nothing was omitted from any retained span, so the package carries no omission record.  No retained line contains a citation, so the wrapper carries no bibliography and no bibliography entry.  No retained line contains a figure environment, an image-inclusion command or drawing code, so no image is delivered and the package declares no asset dependency.  Editorial formatting only, in addition: an AMS \texttt{amsart} preamble that restates the article's own declarations and macros used by the retained lines, namely the environments \texttt{lem} on separate counters, together with the standard packages those lines need.  The retained environments print in this document as Lemma 1 in order of appearance, whereas the article prints the same items with the numbering of its own full document.  The complete native source of the article as distributed in the arXiv e-print for this exact version is bundled unchanged as \texttt{sources/205891/source0.tex}.
\begin{lem}
\label{lem:Stone=000020Dini}Let $K$ be a compact topological space
and $M$ be a family of continuous functions on $K$ with the following
properties:
\begin{enumerate}
\item For any $x\in K$, there exists $u\in M$ with a strict maximum in
$x$.
\item For all $u\in M$, $a\in\mathbb{R}$, and $b>0$, also $a+bu\in M$.
\item For all $u,w\in M$, also $\max(u,w)\in M$.
\end{enumerate}
Then $M$ is dense in $C(K)$.

\end{lem}

\begin{proof}
Let $w\in C(K)$ and $\varepsilon>0$ be given. For $x\in K$, let
$u_{x}\in M$ be a function satisfying $u_{x}(x)=0>u(K\setminus\{x\})$
by (1) and (2). Let $U(x)$ be a neighborhood of $x$ where $\sup_{U}(w)-\inf_{U}(w)\leq\varepsilon$.
Choose $a_{x}>0$ sufficiently large such that 
\[
w_{x}:=a_{x}u_{x}+w(x)\leq w\quad\text{outside }U(x).
\]
Inside $U(x)$, this function is bounded from above by $w(x)\leq w+\varepsilon$,
altogether $w_{x}\leq w+\varepsilon$ in $K$. In $x$, we have the
identity $w_{x}(x)=w(x)$ and $W_{x}:=\{|w-w_{x}|<\varepsilon\}$
is a neighborhood of $x$. Due to compactness, there exist finitely
many points $x_{1},\dots,x_{m}$ such that $W_{x_{1}},\dots,W_{x_{m}}$
cover $K$. Then $w_{\varepsilon}:=\max\{w_{x_{1}},\dots,w_{x_{m}}\}\in M$
by (3) and $w-\varepsilon\leq w_{\varepsilon}\leq w+\varepsilon$.
\end{proof}

\end{document}
